% native TikZ fragment (batch_095, kind=single_function_options)
% 2013 山东文 9 (C)  y=x cos x + sin x
%% shape C : odd; f(0)=0 flat at origin f'(0)=0; decreasing through O; max +3.12 at x~-3.0 (min mirrored); f(pi)=-pi<0
\usetikzlibrary{arrows.meta}
\begin{tikzpicture}[x=1cm,y=1cm,line width=.6pt,>=Stealth,font=\normalsize]
\draw[->] (-3.95,0)--(3.95,0) node[right,inner sep=0pt,xshift=2pt]{$x$};
\draw[->] (0,-3.95)--(0,3.95) node[above,inner sep=0pt,yshift=2pt]{$y$};
\draw (3.1416,-0.10)--(3.1416,0.10);
\node[anchor=north,inner sep=0pt] at (3.1416,-0.26) {$\pi$};
\node[anchor=north west,inner sep=0pt] at (0.30,-0.26) {$O$};
\draw[domain=-3.80:3.80,samples=281] plot (\x,{\x*cos(deg(\x))-sin(deg(\x))});
\end{tikzpicture}
